
\documentclass{article}
\usepackage{colortbl}
\usepackage{makecell}
\usepackage{multirow}
\usepackage{supertabular}

\begin{document}

\newcounter{utterance}

\centering \large Interaction Transcript for game `eqbench', experiment `eq\_bench', episode 160 with reflexion{-}Qwen3.5{-}27B.
\vspace{24pt}

{ \footnotesize  \setcounter{utterance}{1}
\setlength{\tabcolsep}{0pt}
\begin{supertabular}{c@{$\;$}|p{.15\linewidth}@{}p{.15\linewidth}p{.15\linewidth}p{.15\linewidth}p{.15\linewidth}p{.15\linewidth}}
    \# & \multicolumn{2}{c}{Player} && \multicolumn{2}{c}{Game Master} \\
    \hline

    \theutterance \stepcounter{utterance}  
    & & & \multicolumn{4}{p{0.6\linewidth}}{
        \cellcolor[rgb]{0.9,0.9,0.9}{
            \makecell[{{p{\linewidth}}}]{
                \texttt{\tiny{[P1$\langle$GM]}}
                \texttt{You have attempted this same task before and did not succeed. These are the plans you wrote after those attempts.} \\
\texttt{Trial \#1: New plan: My previous attempt failed because I underestimated the intensity of "Regretful" and overestimated "Defensive" without fully analyzing the specific emotional arc of the final lines. Richie's accusation ("That's why you left, isn't it?") followed by Jules's confirmation ("I left because you wouldn't listen") directly challenges Richie's self{-}perception and highlights his failure to communicate, which is a primary driver of regret. While Richie is defensive throughout, the dialogue ends on a note of realization of his own fault rather than just a defensive posture. Next time, I will re{-}evaluate the final exchange to identify the shift from conflict to realization, assigning a higher score to "Regretful" (likely the dominant emotion) and adjusting "Defensive" downward to reflect that the conversation has moved past pure defense into an acknowledgment of past failures. I will also ensure "Shocked" remains low as the conflict was expected, and "Proud" remains near zero as the dialogue strips away any ego.} \\
\\ 
\texttt{Trial \#2: New plan: My previous attempt failed because I misidentified the dominant emotional state at the very end of the dialogue. While Richie is certainly defensive throughout the exchange, the final line from Jules ("I left because you wouldn't listen. You never did.") serves as a definitive, unrefuted indictment of Richie's character. This specific phrasing shifts the emotional weight from active defense to a crushing realization of his own inability to connect, which is the definition of regret. In my next attempt, I will prioritize the final sentence as the emotional anchor. I will assign the highest score to "Regretful" (e.g., 9 or 10) because the dialogue ends with Richie being confronted with his own fatal flaw without a chance to retort. I will lower "Defensive" (e.g., to 3 or 4) to reflect that the conversation has moved past the back{-}and{-}forth arguing into a moment of painful truth. "Shocked" and "Proud" will remain at 0 or 1, as the conflict was anticipated and Richie's ego has been dismantled. The key is to recognize that the silence after Jules's final statement implies Richie is processing his failure, not fighting back.} \\
\\ 
\texttt{Trial \#3: New plan: My previous attempts failed because I misjudged the emotional trajectory of the final exchange. While Richie is defensive throughout the dialogue, the specific phrasing of Jules's final line ("I left because you wouldn't listen. You never did.") acts as a definitive, unrefuted indictment that shifts the emotional weight from active defense to a crushing realization of his own failure. This realization is the core of "Regretful." In my next attempt, I will treat the final sentence as the primary emotional anchor, assigning the highest possible score to "Regretful" (10) to reflect the weight of this unchallenged truth. I will lower "Defensive" (to 3 or 4) to indicate that the conversation has moved past the back{-}and{-}forth arguing into a moment of painful acceptance. "Shocked" and "Proud" will remain at 0 or 1, as the conflict was expected and Richie's ego has been dismantled. The key is to recognize that the silence after Jules's final statement implies Richie is processing his failure, not fighting back.} \\
\\ 
\\ 
\\ 
\texttt{Your task is to predict the likely emotional responses of a character in this dialogue:} \\
\\ 
\texttt{Richie: Jules Baxter, as I live and breathe. Never thought I'd see you at a festival again.} \\
\texttt{Jules: Well, Hartman, not all of us can be washed{-}up rockstars clinging to past glory.} \\
\texttt{Richie: Ouch. That's a bit harsh, don't you think?} \\
\texttt{Jules: You always did have a knack for understatement.} \\
\texttt{Richie: And you always had a knack for being a buzzkill. } \\
\texttt{Jules: I was the only one who ever took the music seriously, Richie. You were too busy playing the rockstar.} \\
\texttt{Richie: And you were too busy playing the martyr. } \\
\texttt{Jules: I was trying to create something meaningful. You just wanted the fame.} \\
\texttt{Richie: And you just wanted to control everything. That's why you left, isn't it?} \\
\texttt{Jules: I left because you wouldn’t listen. You never did.} \\
\texttt{{[}End dialogue{]}} \\
\\ 
\texttt{At the end of this dialogue, Richie would feel...} \\
\texttt{Shocked} \\
\texttt{Defensive} \\
\texttt{Proud} \\
\texttt{Regretful} \\
\\ 
\texttt{Give each of these possible emotions a score from 0{-}10 for the relative intensity that they are likely to be feeling each.} \\
\\ 
\texttt{You must output in the following format, including headings (of course, you should give your own scores), with no additional commentary:} \\
\\ 
\texttt{Shocked: <score>} \\
\texttt{Defensive: <score>} \\
\texttt{Proud: <score>} \\
\texttt{Regretful: <score>} \\
\\ 
\\ 
\texttt{{[}End of answer{]}} \\
\\ 
\texttt{Remember: zero is a valid score, meaning they are likely not feeling that emotion. You must score at least one emotion > 0.} \\
\\ 
\texttt{Your answer:} \\
            }
        }
    }
     \\ \\

    \theutterance \stepcounter{utterance}  
    & \multicolumn{4}{p{0.6\linewidth}}{
        \cellcolor[rgb]{0.561,0.737,0.561}{
            \makecell[{{p{\linewidth}}}]{
                \texttt{\tiny{[P1$\rangle$GM]}}
                \texttt{Shocked: 1} \\
\texttt{Defensive: 4} \\
\texttt{Proud: 0} \\
\texttt{Regretful: 9} \\
            }
        }
    }
    & & \\ \\

    \theutterance \stepcounter{utterance}  
    & & & \multicolumn{2}{p{0.3\linewidth}}{
        \cellcolor[rgb]{0.9,0.9,0.9}{
            \makecell[{{p{\linewidth}}}]{
                \texttt{\tiny{[GM$|$GM]}}
                \texttt{Shocked: 1} \\
\texttt{Defensive: 4} \\
\texttt{Proud: 0} \\
\texttt{Regretful: 9} \\
            }
        }
    }
    & & \\ \\

    \theutterance \stepcounter{utterance}  
    & & & \multicolumn{2}{p{0.3\linewidth}}{
        \cellcolor[rgb]{0.9,0.9,0.9}{
            \makecell[{{p{\linewidth}}}]{
                \texttt{\tiny{[GM$|$GM]}}
                \texttt{Shocked: 2} \\
\texttt{Defensive: 8} \\
\texttt{Proud: 4} \\
\texttt{Regretful: 6} \\
            }
        }
    }
    & & \\ \\

    \theutterance \stepcounter{utterance}  
    & & & \multicolumn{2}{p{0.3\linewidth}}{
        \cellcolor[rgb]{0.9,0.9,0.9}{
            \makecell[{{p{\linewidth}}}]{
                \texttt{\tiny{[GM$|$GM]}}
                \texttt{game\_result = LOSE} \\
            }
        }
    }
    & & \\ \\

    \theutterance \stepcounter{utterance}  
    & & \multicolumn{4}{p{0.6\linewidth}}{
        \cellcolor[rgb]{0.95,0.95,0.95}{
            \makecell[{{p{\linewidth}}}]{
                \texttt{\tiny{[P1$\leftrightarrow$P]}}
                \texttt{Post{-}episode Reflexion (not sent to GM):} \\
\texttt{New plan: My previous attempts failed because I incorrectly identified "Regretful" as the dominant emotion, likely misinterpreting the silence after Jules's final line as acceptance rather than a defensive shutdown. Richie's entire dialogue is characterized by deflection and counter{-}accusations (e.g., "playing the martyr," "wanted to control everything"). When Jules delivers the final blow ("You never did"), Richie's established pattern suggests he would double down on his defensiveness rather than immediately feel regret. The "Regretful" score was too high, and "Defensive" was too low. Next time, I will prioritize the character's consistent behavioral pattern of deflection over the theoretical implication of the final line. I will assign the highest score to "Defensive" (e.g., 9 or 10) to reflect his refusal to accept the accusation, lower "Regretful" significantly (e.g., 2 or 3) as he is not yet admitting fault, and keep "Shocked" and "Proud" at 0 or 1. The key is to recognize that Richie's character arc in this snippet is one of resistance, not realization.} \\
            }
        }
    }
    & \\ \\

\end{supertabular}
}

\end{document}
